% =============================================================================
%  Notas de clase — Sesión 01: Presentación del Curso
%  Programación Avanzada — Universidad Panamericana
%  Compilar: pdflatex sesion01_presentacion.tex
% =============================================================================
\documentclass[12pt, oneside]{article}
\usepackage{progavanzada}
\usepackage{multicol}

\begin{document}

\portada{Sesión 01}{Presentación del Curso}{2026}

\tableofcontents
\newpage

% =============================================================================
\section{Bienvenida}
% =============================================================================

Bienvenido a \textbf{Programación Avanzada}. Este curso es una inmersión
profunda en C++, uno de los lenguajes más relevantes y demandantes de la
industria tecnológica. Aquí no solo aprenderás sintaxis: entenderás
\textit{por qué} el computador hace lo que hace, y tendrás control directo
sobre ello.

\begin{recuerda}
  Lo que distingue a este curso de otros de programación es el énfasis en
  \textbf{entender la máquina}. Aprenderás a gestionar memoria, a diseñar
  funciones eficientes y a pensar como lo haría el compilador.
\end{recuerda}

% =============================================================================
\section{Syllabus}
% =============================================================================

\subsection{Unidad 1 — Fundamentos del lenguaje C++}
\begin{itemize}
  \item Introducción a C++: historia, compiladores y entornos de trabajo
  \item Tipos de datos, variables y operadores
  \item Estructuras de control: \texttt{if}, \texttt{switch}, \texttt{while},
        \texttt{for}
  \item Funciones: declaración, definición, ámbito y recursión
\end{itemize}

\subsection{Unidad 2 — Memoria y apuntadores}
\begin{itemize}
  \item Modelo de memoria: stack y heap
  \item Apuntadores y referencias
  \item Paso por valor, por referencia y por apuntador
  \item Arreglos y aritmética de apuntadores
  \item Memoria dinámica: \texttt{new} y \texttt{delete}
\end{itemize}

\subsection{Unidad 3 — Programación orientada a objetos}
\begin{itemize}
  \item Clases y objetos
  \item Encapsulamiento, herencia y polimorfismo
  \item Constructores, destructores y sobrecarga de operadores
  \item Clases abstractas e interfaces
\end{itemize}

\subsection{Unidad 4 — Estructuras de datos y algoritmos}
\begin{itemize}
  \item Algoritmos de ordenamiento: Bubble Sort, Quick Sort, Merge Sort
  \item Búsqueda lineal y binaria
  \item Introducción a la STL: \texttt{vector}, \texttt{map}, \texttt{set}
\end{itemize}

% =============================================================================
\section{Cómo se califica}
% =============================================================================

La calificación final se compone de tres partes iguales:

\begin{center}
\begin{tabular}{clc}
  \toprule
  \textbf{Porcentaje} & \textbf{Rubro} & \textbf{Descripción} \\
  \midrule
  33\% & Tareas y trabajos en clase
        & Ejercicios, cuadernos y proyectos entregados \\
  33\% & Asistencia y participación
        & Presencia, contribución en discusiones y actividades \\
  33\% & Examen
        & Evaluación escrita y/o práctica por unidad \\
  \bottomrule
\end{tabular}
\end{center}

\begin{consejo}
  \textbf{La participación vale tanto como el examen.} Las preguntas en clase,
  los comentarios en las discusiones y tu presencia activa construyen el 33\%
  de tu calificación desde el primer día. No hay forma de recuperar esos puntos
  al final del semestre.
\end{consejo}

\begin{advertencia}
  Las tareas se entregan en la fecha acordada. Cada día de retraso puede
  reducir la calificación de ese entregable. Consulta la política específica
  de tu profesor al inicio del semestre.
\end{advertencia}

% =============================================================================
\section{Bibliografía recomendada}
% =============================================================================

\subsection{Libros principales}

\textbf{1. The C++ Programming Language} — Bjarne Stroustrup\\
4.ª edición, Addison-Wesley, 2013.\\
El libro definitivo del creador del lenguaje. Denso y técnico, pero
la referencia más completa que existe. Ideal para consultar cuando quieras
entender \textit{por qué} algo funciona así en C++.\\
\url{https://www.stroustrup.com/4th.html}

\vspace{0.8em}
\textbf{2. A Tour of C++} — Bjarne Stroustrup\\
3.ª edición, Addison-Wesley, 2022.\\
Versión condensada del anterior. Recorre el lenguaje completo en $\sim$300
páginas. Excelente para tener una visión panorámica rápida.\\
\url{https://www.stroustrup.com/tour3.html}

\vspace{0.8em}
\textbf{3. C++ Primer} — Lippman, Lajoie, Moo\\
5.ª edición, Addison-Wesley, 2012.\\
El libro introductorio más utilizado en cursos universitarios de C++.
Paso a paso, con muchos ejemplos y ejercicios. Muy recomendado si es tu
primer contacto formal con el lenguaje.

\vspace{0.8em}
\textbf{4. The C Programming Language} — Kernighan \& Ritchie\\
2.ª edición, Prentice Hall, 1988.\\
El legendario ``K\&R''. C++ nació de C, y entender C ilumina muchas
decisiones de diseño de C++. Clásico imprescindible para cualquier
programador serio.

\subsection{Recursos en línea (gratuitos)}

\begin{itemize}
  \item \textbf{cppreference.com} — \url{https://en.cppreference.com}\\
        La referencia técnica más completa y actualizada del lenguaje y la STL.

  \item \textbf{learncpp.com} — \url{https://www.learncpp.com}\\
        Tutorial gratuito muy completo, con ejercicios. Excelente complemento
        a las clases.

  \item \textbf{C++ Core Guidelines} — \url{https://isocpp.github.io/CppCoreGuidelines}\\
        Guía de buenas prácticas escrita por Stroustrup y Herb Sutter.
        Aprende no solo a que el código funcione, sino a que sea correcto.

  \item \textbf{Compiler Explorer (Godbolt)} — \url{https://godbolt.org}\\
        Herramienta online que muestra el ensamblador generado por tu código.
        Fascinante para entender qué hace realmente el compilador.
\end{itemize}

\begin{tecnico}
  \textbf{¿Por qué C++ sigue siendo relevante en 2026?}\\
  Según el índice TIOBE, C++ se mantiene consistentemente entre los 3
  lenguajes más usados del mundo. Es el lenguaje dominante en
  videojuegos (Unreal Engine), sistemas embebidos, motores de bases de
  datos, navegadores web y software de alto rendimiento. El mercado
  laboral para programadores C++ ofrece algunos de los salarios más
  altos en la industria.\\
  \textit{Referencia:} \url{https://www.tiobe.com/tiobe-index/cplusplus/}
\end{tecnico}

% =============================================================================
\section{Actividades de la sesión}
% =============================================================================

\subsection{Padlet — ¿Qué emoción te produce programar?}

En el Padlet del curso encontrarás un tablero compartido.
Tu tarea es \textbf{agregar un post-it} con un meme, GIF o imagen que
represente cómo te sientes cuando programas.

No hay respuesta correcta o incorrecta. La intención es romper el hielo
y conocernos como grupo.

\begin{center}
  \large\textbf{Padlet del curso:}\\[6pt]
  \url{https://padlet.com}\\[4pt]
  {\small (El enlace exacto lo comparte el profesor al inicio de clase)}
\end{center}

\begin{consejo}
  Los mejores programadores no son los que no sienten frustración —
  son los que aprendieron a convertir la frustración en curiosidad.
  Si tu meme es de frustración total, bienvenido: estás en el lugar correcto.
\end{consejo}

\subsection{Lluvia de ideas — ¿Qué se necesita para ser un buen programador hoy?}

Reflexiona sobre estas preguntas antes de compartir tus ideas con el grupo:

\begin{itemize}
  \item ¿Qué habilidades técnicas consideras esenciales?
  \item ¿Qué habilidades \textit{no técnicas} crees que importan más de lo que
        la gente piensa?
  \item ¿Cómo cambia la respuesta si consideramos herramientas de IA como
        GitHub Copilot, ChatGPT o Claude?
  \item ¿Qué diferencia a alguien que \textit{usa} un lenguaje de alguien que
        realmente lo \textit{domina}?
\end{itemize}

\vspace{0.5cm}
\noindent\textbf{Mis notas de la discusión:}

\vspace{3cm}
\noindent\rule{\textwidth}{0.4pt}
\vspace{2.5cm}
\noindent\rule{\textwidth}{0.4pt}
\vspace{2.5cm}
\noindent\rule{\textwidth}{0.4pt}

\subsection{Ejercicio de mecanografía — 10FastFingers}

Uno de los factores que más impacta tu flujo de trabajo como programador es la
\textbf{velocidad y precisión al escribir}. No se trata de escribir rápido a
expensas de pensar — se trata de que tu velocidad de escritura no sea el
cuello de botella cuando ya tienes la solución clara en la cabeza.

\begin{center}
  \large\textbf{Plataforma:}\\[6pt]
  \url{https://10fastfingers.com}\\[4pt]
  {\small Modo: Español — Test de 1 minuto}
\end{center}

\noindent\textbf{Registra tu resultado aquí:}

\begin{center}
\begin{tabular}{lc}
  \toprule
  \textbf{Métrica} & \textbf{Mi resultado} \\
  \midrule
  Palabras por minuto (PPM) & \hspace{3cm} \\[8pt]
  Precisión (\%)            & \hspace{3cm} \\[8pt]
  Fecha                     & \hspace{3cm} \\
  \bottomrule
\end{tabular}
\end{center}

\begin{tecnico}
  \textbf{¿Cuánto es normal?}\\
  El promedio de un adulto es $\sim$40 PPM. Un programador experimentado
  suele estar entre 60--80 PPM. Mecanógrafos profesionales superan los
  100 PPM. La mayoría de las personas mejoran 10--15 PPM en pocas semanas
  de práctica deliberada. Lo importante no es el número de hoy, sino la
  tendencia al final del semestre.\\
  \textit{Referencia:} Dhakal et al. (2018). ``Observations on Typing from
  136 Million Keystrokes''. \textit{CHI 2018}. \url{https://doi.org/10.1145/3173574.3174220}
\end{tecnico}

% =============================================================================
\section{Lo que viene}
% =============================================================================

En la \textbf{Sesión 02} comenzamos con C++ directamente:

\begin{itemize}
  \item Instalación y verificación de \texttt{g++} en tu equipo
  \item Primer programa: \texttt{Hello World}
  \item Compilar y ejecutar desde la consola
  \item Apertura del cuaderno interactivo en Jupyter con xeus-cling
\end{itemize}

\begin{recuerda}
  Para la próxima sesión necesitas tener acceso a una terminal con
  \texttt{g++} instalado \textbf{o} acceso a
  \href{https://shell.cloud.google.com}{Google Cloud Shell}
  (no requiere instalación, funciona en cualquier navegador).
\end{recuerda}

\end{document}
